% ============================================================
% D03 — Arsitektur dan Aktor
% ============================================================
\flowmeta{D03}{Arsitektur dan Aktor}{Umum}

\begin{flowsection}{Ringkasan Arsitektur}
TUSKO dibangun dengan arsitektur \emph{decoupled}: antarmuka pengguna (frontend)
dan layanan data (backend API) berjalan terpisah dan berkomunikasi melalui
permintaan HTTP/JSON. Gambaran lapisan sistem:
\begin{itemize}
  \item \textbf{Frontend} --- aplikasi web React (SPA) untuk storefront pelanggan
        dan panel admin ERP.
  \item \textbf{API Backend} --- layanan REST berbasis Laravel dengan autentikasi
        token (Sanctum) dan login Google.
  \item \textbf{Basis Data} --- basis data relasional (MySQL) untuk seluruh data
        pengguna, produk, pesanan, inventori, dan keuangan.
  \item \textbf{Layanan Pihak Ketiga} --- Midtrans (pembayaran), Biteship
        (logistik/kurir), penyimpanan berkas (R2/S3), dan pengiriman email.
\end{itemize}
\end{flowsection}

\shot{gambar/umum/d03-1.png}{Diagram arsitektur sistem TUSKO beserta aktor dan layanan pihak ketiga}{fig:d03-1}

\begin{flowsection}{Aktor Sistem}
\begin{itemize}
  \item \textbf{Pelanggan} --- pengguna terdaftar yang berbelanja; dapat melihat
        katalog, mengelola akun, dan melacak pesanan.
  \item \textbf{Tamu} --- pengunjung yang dapat melihat katalog publik tanpa login.
  \item \textbf{Admin} --- pengelola toko yang mengakses panel ERP sesuai peran
        (menu ditentukan oleh peran masing-masing).
\end{itemize}
\end{flowsection}

\begin{flowsection}{Aliran Data Utama}
Pelanggan/tamu berinteraksi dengan frontend. Frontend mengirim permintaan ke API
backend; backend memvalidasi, memproses aturan bisnis, menyimpan/membaca dari
basis data, dan bila perlu memanggil layanan pihak ketiga (mis. membuat
pembayaran ke Midtrans atau pengiriman ke Biteship). Hasilnya dikembalikan ke
frontend untuk ditampilkan.
\end{flowsection}

\begin{flowsection}{Kaitan Modul}
Frontend: \texttt{frontend/src/} (React, Tailwind, Lucide). Backend:
\texttt{backend/app/} (Laravel, controller/service). Rute: \texttt{backend/routes/api.php}.
\end{flowsection}

\docver{v1.0}{Sep 2026}
